% !TEX root = ../../mixed-sklyanin-cubics.tex
% Current result: 2026-11-22 / v62.
\subsection{Clifford 九変数族の幾何的接類と局所比較}
\label{ifk:subsec:clifford-application}

中心変数を加えない九変数族は三次四様体の次元に属する。
本節では実相対模型と原点の rotation 付き核を用い、
主稿の rotation を持つ変形の形式認識を適用する。
まず残余圏の九つの接方向に非幾何的成分がないことを確認し、
次に原点近傍の実際の可換比較へ進む。

定理\ref{ifc:thm:family}の $m=0$ の族を取り、
$0\in\operatorname{Spec}R$ を $\Lambda=0$ とする。
原点の商と残余圏を
\[
 A_0=(S_1\otimes S_2)/(a_1\Omega_{1,0}+a_2\Omega_{2,0}),
 \qquad \mathcal K_0=\Ku_{\mathrm{coh}}(A_0)
\]
と書く。
各因子の点曲線の非零方程式を $f_i$ とし、
$X_0=V(a_1f_1+a_2f_2)\subset\mathbb P^5$ と置く。
補題\ref{ebk:lem:boundary}は、この $\mathcal K_0$ と $\Ku(X_0)$ の間に
実際の dg Morita 同値 $E$ と
\begin{equation}\label{ifk:eq:marked-origin}
 E\rho_0\simeq\rho_{X_0}E
\end{equation}
を、指定三乗周期とともに与える。
同補題は二捩れ因子を許すので $S_2=S(1,1,\tau)$ に適用できる。
この入力は単なる非標識 Morita 同値より強い。

\begin{proposition}[九方向の非幾何的成分の消滅]
\label{ifk:prop:clifford-geometric-ks}
定理\ref{iff:thm:smooth-neighborhood}の実相対模型に対し、
原点における残余圏の Kodaira--Spencer 写像は
\begin{equation}\label{ifk:eq:clifford-ks}
 \operatorname{Mat}_3(\C)
   \xrightarrow{\ \operatorname{KS}_{\mathcal K}\ }
 \mathrm{HH}^2(\mathcal K_0)
   \xrightarrow{\ E_*\ }
 \mathrm{HH}^2(\Ku(X_0))
\end{equation}
の像を $\operatorname{im}\eta_{X_0}\simeq(J_{a_1f_1+a_2f_2})_3$ に持つ。
特に、二つの非恒等 rotation 固有空間への射影は
九方向全てについて零である。
\end{proposition}

\begin{proof}
定理\ref{iff:thm:smooth-neighborhood}により、
原点を含む開集合上に実相対 smooth proper 模型と
内部 shift から来る自己同値 $\rho$ がある。
その Artin 底変更により、任意の
$\Lambda_{ab}=\epsilon M_{ab}$ に対して
$\mathcal K_0$ の実際の Morita 変形と $\rho_0$ の持上げを得る。
任意の $M$ に補題\paperref{main}{ifk:lem:lifted-autoequivalence}を適用すると、
$\operatorname{KS}_{\mathcal K}(M)$ は $\rho_0$ に固定される。
式\eqref{ifk:eq:marked-origin}は、実際の可逆双加群による共役作用を
intertwine するので、その $E_*$ 像は $\rho_{X_0}$ に固定される。
命題\paperref{main}{ifk:prop:geometric-fixed-part}から結論を得る。
\end{proof}

\begin{corollary}[四次元九変数族の原点近傍における可換比較]
\label{ifk:cor:clifford-fourfold-neighborhood}
$m=0$ の九変数族について、原点を含む Zariski 開集合
$V\subset\mathbb A^9_\Lambda$ が存在し、
各複素点の全 coherent Kuznetsov 成分は、
滑らかな可換三次四様体の Kuznetsov 成分に dg Morita 同値である。

具体的には、式\eqref{ifk:eq:marked-origin}の実際の核 $E$ を固定すると、
原点の \'etale 近傍 $V'\to\mathbb A^9_\Lambda$ 上に、
$X_0$ を延長する滑らかな三次四様体族と、
$E$ を延長する perfect な実 Morita 比較核を構成できる。
この同値は任意の導来底変更に保たれる。
\end{corollary}

\begin{proof}
定理\ref{iff:thm:smooth-neighborhood}は、原点を含む実開集合上に
相対 smooth proper な残余圏族と、その実際の相対自己同値 $\rho$ を与える。
補題\ref{ebk:lem:boundary}の核 $E$ は
式\eqref{ifk:eq:marked-origin}を満たす。
従って主稿の定理\paperref{main}{ifk:thm:rotation-neighborhood}の
全ての仮定が満たされ、その結論を適用できる。
複素ファイバーで残余模型が全 coherent 残余圏に一致することも、
定理\ref{iff:thm:smooth-neighborhood}の結論である。
\end{proof}

\begin{remark}[接類と局所比較の範囲]
\label{ifk:rem:tangent-scope}
命題\ref{ifc:prop:tangent}は周囲代数の
$\mathrm{HH}^2(S_1\otimes S_2)_{(-2,2)}$ における九方向の独立性を証明する。
命題\ref{ifk:prop:clifford-geometric-ks}は残余圏への像の位置を制限するものであり、
写像\eqref{ifk:eq:clifford-ks}の単射性や rank 九を証明するものではない。
実相対残余模型は代数の次数別自由性だけから得たのではなく、
定理\ref{iff:thm:smooth-neighborhood}の相対 perfect 対角と
実自己同値の構成を用いている。

系\ref{ifk:cor:clifford-fourfold-neighborhood}では、一次成分の消滅だけを
全次数の比較へ置き換えたのではない。
主稿の定理\paperref{main}{ifk:thm:rotation-formal}による
有限位数作用の固定形式部分の認識と、
比較核スタックの有限表示を合わせて、実際の近傍の核を得ている。
全ての smooth 係数への延長、境界での比較核の保持、
および核による全 rotation・周期・中心作用の同時保存は、依然として別の課題である。
\end{remark}
